\section{Conclusion and Future Work}
\label{sec:conclusion}

\paragraph{This is a first step.}
This paper does not validate the OpenBook QR Score, and its scores should not be read as measurements of how good or how reusable the ten textbooks are. It is a fast, low-cost first step: an argument for what an objective OER score must be, a specification of one candidate, and an illustration that applies the candidate to ten books in minutes of wall-clock time per pass. Every score comes from AI agents of one model family, from three sampled chapters per book, and no human reviewed any resource.

\paragraph{What the first step showed.}
We argued that deciding which open educational resources to adopt needs a shared instrument that is anchored, two-part, reproducible and validated, and that existing rubrics, as surveyed in the literature, meet these requirements only partly. We proposed the OpenBook QR Score, a ten-metric, 100-point instrument whose report keeps quality and reusability visible as QR 76 (Q 37 / R 39) and whose decision rule guards against a high total hiding a failed metric. The illustration showed the instrument's value (it points to concrete fixes, and two passes agreed within one point on 88 percent of metric scores) and its weaknesses (descriptors for restrictive licenses are loose, practice with feedback is unstable across chapters, and quality is judged mainly from structure, not from checked content).

\paragraph{The next step.}
The next step, for a future paper, is to verify the content of sample chapters with a low-cost model, Claude Haiku, in the evidence-grounded, repeated and calibrated pipeline of Section~\ref{sec:nextstep}. It targets the weakest part of this paper, the accuracy of the quality metrics, and it will be tested first on chapters seeded with known errors.

\paragraph{The steps after that.}
Four further steps follow. First, tighten the descriptors where the illustration exposed drift, starting with licenses and learning outcomes. Second, run the human and agent agreement studies of Section~\ref{sec:validation}, with at least two independent human reviewers per resource and several chapters per book. Third, test whether scores predict learner outcomes in controlled comparisons of an original resource and an improved version. Fourth, publish scores with their evidence and version in an open catalog, so that the community can dispute and improve them. Until those steps are done, every OpenBook QR Score should be labeled provisional.
